\section{Discussion}\label{sec:discussion}
\subsection{Generalization to More Types of Terrain}
The current framework relies on manual terrain abstraction and classification, which requires expert knowledge to define terrain types and assign them to discrete symbolic categories. This process typically involves analyzing physical features such as terrain shape, height variation, and support area to determine whether a terrain segment should be labeled as flat, sparse, dense, high, gap, etc. While effective, this expert-driven approach may limit scalability and generalization to novel environments or terrain conditions. Future extensions could incorporate data-driven terrain classification and generation methods, such as supervised learning or clustering based on 3D point cloud statistics to automate the terrain abstraction process and reduce reliance on human heuristics.

\subsection{Optimality}
The synthesis process focuses on guaranteeing feasibility and correctness of transitions but does not account for the relative cost or quality of different feasible gaits. When multiple locomotion gaits (e.g., trotting, bounding, leaping) are viable for a given symbolic transition, the current method does not evaluate their optimality in terms of energy efficiency, robustness, execution time, or other performance metrics. Instead, it selects the first gait that satisfies the physical feasibility conditions via MICP. Integrating cost-aware synthesis would enable the framework to prioritize higher-quality behaviors and make more informed decisions under trade-offs.

\subsection{Perception Module}
The proposed framework is modular and can be directly integrated with a wide range of perception pipelines that provide terrain geometry in the form of segmented polygonal regions \citep{miki2022elevation}. Existing perception modules that perform terrain segmentation using LiDAR, RGB-D sensors, or stereo vision can be adapted to supply the required input to the MICP planner and symbolic abstraction layers. \revised{In this work we deliberately treated the perception layer as out of scope and used motion-capture-aided manual labeling on hardware (Sec.~\ref{sec:hardware}); a fully autonomous deployment will need to address how perception uncertainty---segmentation noise, polygon misalignment, and per-frame perception delay---propagates through the MICP feasibility certificates and the symbolic-level realizability guarantees. In the current framework, runtime symbolic repair already provides a partial fallback: when perceived terrain causes the symbolic specification to become unrealizable, repair triggers and synthesizes new skills. Quantifying the resilience of this fallback under realistic perception noise is an important next step.}

\revised{\subsection{Reactivity and Update Frequencies}
Across the experiments reported in this article, the symbolic strategy automaton evolves at the timescale of one symbolic transition, which corresponds to one online MICP solve---typically hundreds of milliseconds to a few seconds (Tables~\ref{tab:micp_speed} and~\ref{tab:micp_speed_hardware}). Two complementary directions exist for tightening this loop in future work.

First, more frequent MICP updates---e.g., re-solving the online gait-fixed MICP at every NMPC tick rather than once per symbolic transition---are easily adaptable in the proposed framework: the delay-aware coordination scheme already supports asynchronous MICP solves of arbitrary duration, so any speed-up of the underlying mixed-integer solver, or any per-transition simplification, translates directly into a higher MICP update rate. The current limit is therefore the MICP solve time itself rather than any architectural constraint.

Second, more frequent symbolic-level planning is also achievable but requires careful design of the symbolic abstraction. Higher-frequency symbolic decisions would generally call for a finer abstraction---more atomic propositions per cell, finer terrain or robot-state granularity, or a larger skill repertoire. Reactive synthesis, however, scales exponentially in the number of variables in the specification~\citep{bloem2012synthesis}, so any finer abstraction must be balanced against synthesis cost. The high-level manager and symbolic repair adopted in this article already mitigate this scalability barrier; pushing toward truly real-time symbolic planning is an interesting future direction at the abstraction level rather than a limit of the framework itself.}

\revised{\subsection{Stronger Baseline Comparisons as Future Work}
The experimental validation in this article compares against a pure-MIP planner and a nearest-feasible-foothold heuristic. Two stronger comparison points that we did not implement here remain valuable future targets: (i) optimization-based planners such as TOWR-style phase-based end-effector parameterization~\citep{winkler2018gait}, which jointly optimize gait timing and footholds in continuous time; and (ii) end-to-end perceptive reinforcement-learning controllers, in particular parkour-style policies such as~\citep{zhuang2023robot,hoeller2023anymal,cheng2024extreme,luo2024pie}. A fair comparison against these two methods requires either equipping them with an additional discrete terrain selection mechanism or a comparable high-level navigation layer, which we did not pursue in the present validation.}

\revised{\subsection{Extension to Bipedal and Other Platforms}
This article focuses on quadrupedal locomotion. The framework's symbolic-level machinery---terrain abstraction, GR(1) specifications, symbolic manager, symbolic repair, runtime repair---is platform-agnostic and applicable directly to other systems such as bipedal robots. The MICP layer requires more careful adaptation: angular-momentum modulation becomes critical even at moderate speeds---a regime where the convex-dynamics simplification adopted here (Eq.~(\ref{eq:dynamics})) becomes a stronger limitation. Generalizing to humanoids would also require revisiting the kinematic-feasibility re-targeting MICP to account for arm-swing and balance constraints. We see this as a future direction.}
